\subsection{电化学阻抗谱}

稳态极化曲线把所有步骤的过电位加总成一个数；阻抗谱则用一串不同频率的小幅正弦扰动，把总响应按\textbf{时间尺度}拆开——快步骤（溶液电阻、双电层充电）在高频段显现，慢步骤（电荷转移、吸附、扩散）在低频段显现\src{S2 §11.5，书内 p.866 / PDF p.354}。分段能力是它最大的价值，也是它最大的风险来源：把频率轴上某一段的形状直接翻译成一个元件，很容易得到「拟合优度很高、物理归属却错」的电路。

本分片源文本自书内 p.866（PDF p.354）起，只包含 §11.5 的后半（应用与读图）与 §11.6 全节；阻抗的基本定义、等效电路推导与 Nyquist／Bode 图的定义在源书 §11.5.1，不在本分片的源文本范围内——它们在前一节（§11.5 阻抗谱基本原理，依源文本写成，标有 \src{} 依据）。下文凡标 \tocheck{} 处，即按学科标准形式重建、未能与源文本对齐的部分。

\begin{probout}
 已知：电极过程的总电流至少含溶液电阻、双电层充电、电荷转移、吸附中间体与扩散五类贡献，稳态方法只能给出它们的总和。
 要问：能否把总响应分解成若干时间常数，分别读出溶液电阻 $R_{\mathrm{s}}$、双电层电容 $C_{\mathrm{dl}}$、电荷转移电阻 $R_{\mathrm{ct}}$ 与扩散阻抗，并据此判断哪一步是速率控制步骤？
 要给出：覆盖足够宽频率范围的一组阻抗数据、一个（或几个）能同时拟合全部频率的等效电路，以及每个参数的物理归属与成立条件——并明确「能拟合」不等于「唯一」。
\end{probout}

\subsubsection{读图实例：甲醇氧化的阻抗谱随电位的变化}

上一节（源书节号 §11.5.1，本项目中即紧接在前的一节）已给出阻抗的定义 $Z(\omega)=E(\omega)/I(\omega)$、Nyquist／Bode 两种画法、等效元件的来源与四类电极的特征谱，以及「高频截距」「半圆直径」「低频 45°斜线」三条读数规则各自成立的条件；本节不再重复，只处理源文本 §11.5.2 的真实数据。

源文本用同一组数据的两套图讨论甲醇在铂系催化剂上的氧化：Nyquist 形式的图 11-42 与以相角为主的图 11-43\src{S2 §11.5.2，书内 p.866 / PDF p.354}。下面把\textbf{观察}与\textbf{推断}严格分开，因为这里的每一步推断都依赖一个前提。

\begin{obs}
随直流电位升高，阻抗谱按四个区间改变形状：
（1）$0.10\sim0.20$ V：谱几乎是一条沿虚轴、相角接近 $90$° 的直线；
（2）$0.20\sim0.33$ V：相角迅速下降，但仍在 $0$° 以上，低频区近似电阻行为；谱线略偏离虚轴；
（3）$0.35$ V 起：谱延伸进第四象限，出现负相角（电感现象）；在 $0.35\sim0.48$ V 区间，电感行为位于低频区，且各条曲线随电位升高整体向高频方向移动；
（4）$0.48$ V 以上：谱突变，分裂为两段分离的曲线，高频段相角越出 $90$°。
但图中纵横坐标与相角的具体数值需对照原书图 11-42、图 11-43\tocheck{未对照原页面图像：源文本对 0.48 V 以上的描述两次出现「超过 90°」（高频段与低频段），疑为 OCR 重复或原文指低频段相角亦越出 0—90° 区间；图的纵轴定义（$-Z''$ 或相角／°）与频率范围无法在本分片源文本内确认}。
\end{obs}

\begin{infer}
（1）接近虚轴且相角接近 $90$° 的直线是理想极化电极（纯电容）特征：在这个电位区间几乎没有法拉第电流，电极／溶液界面只表现为双电层充电。
（2）相角下降而不越过 $0$°，说明出现了耗散性的法拉第通道；谱线开始偏离虚轴，与甲醇离解吸附的开始相对应。
（3）第四象限的负相角是电感行为，源文本把它归属于强吸附中间体 \ce{CO_{ad}} 的吸附弛豫。判据是\textbf{出现电位区间}与已知的 CO 吸附电位区间一致，而不是单看形状——仪器电感也会给第四象限，见下文误判清单第（1）条「第四象限的电感弧」。
（4）源文本采用的判断规则是：\textbf{低频区的过程是速率控制步骤}；因此随电位升高曲线向高频方向移动，就意味着该低频过程的速率在加快。据此，$0.35$ V 以前低频区为电阻行为、速控步是甲醇的离解吸附；$0.35\sim0.48$ V 低频区为电感行为、速控步是生成 \ce{CO_{ad}} 的步骤；$0.48$ V 以上低频区转为钝化行为、速控步是 \ce{CO_{ad} -> CO2}。谱在 $0.48$ V 的突变因此被解释为「CO 的生成与消耗两个步骤中，速控步发生了转移」。
（5）与上一条互补的判据是\textbf{斜率的数值}：$-Z''$ 对 $Z'$ 的斜率明显小于 $45$°（例如落在 $20$°～$30$°）时，通常意味着扩散层已「看到」边界、不再满足半无限扩散的前提，或有吸附弛豫叠加在扩散尾上——此时不能把该段直接当 Warburg 处理。这条判据只在频率范围够低、确实取到了扩散段之后才有意义。
\end{infer}

\begin{dataflow}
从一组阻抗数据到「速控步是谁」的完整链条中，\textbf{检验}（正反扫重合、Kramers--Kronig 残差）与\textbf{数时间常数}（Bode 相角图数峰）两步已写在上一节（源书 §11.5.1），此处不重复；真实数据特有的三步是：
（1）\textbf{定元件归属}——高频截距给 $R_{\mathrm{s}}$，最低频的极限行为决定是扩散（45°）、吸附（低频电容／电感弧）还是钝化（低频大电阻）。
（2）\textbf{跨电位比较}——只有同一套电路、同一套初值在全电位区间内都收敛，才能用「曲线向高频移动」这类趋势论证速率变化。
（3）\textbf{互查}——用独立实验（Tafel 斜率、变转速、原位谱学的覆盖度）验证归属，清单见下文；其中 Tafel 斜率与变转速两条属通用判据，已在上一节的互查清单里连同判据给出。
\end{dataflow}

\subsubsection{拟合的最大风险：参数相关性，以及「拟合好 $\neq$ 电路唯一」}

等效电路拟合的第一风险是\textbf{拟合优度高不等于电路唯一}：不同的电路可以对同一组数据给出几乎相同的拟合曲线。因此逐条列出误判清单仍有必要，尤其是下面几条在真实数据里最容易发生。判据、公式与量化例子见上一节的「误判」清单（十条），此处只列真实数据读图中另有表现的三条：

\paragraph{第四象限的电感弧}引线电感与强吸附中间体（如 \ce{CO_{ad}}）的吸附弛豫都给第四象限的弧。判别办法是改变接线方式或电极几何：\textbf{仪器电感应与频率无关、且能靠改接线消掉}；化学吸附弛豫则出现在确定的电位区间（如读图实例中的 $0.35\sim0.48$ V）。
\paragraph{吸附弛豫与双电层的重叠}吸附中间体的 $R_{\mathrm{ad}}$、$C_{\mathrm{ad}}$ 与 $C_{\mathrm{dl}}$ 的时间常数接近时，二者在 Nyquist 图上不可分辨，只能合并成一个压扁的弧（如读图实例中的甲醇体系）。此时单靠阻抗谱无法把覆盖度信息与双电层信息分开，必须引入原位谱学的覆盖度数据。
\paragraph{薄层与原位谱学联用时的第三套条件}若阻抗数据要与同一电位点的原位红外／拉曼对照，薄层电解池把传质条件改掉了，阻抗谱给出的扩散尾与厚层体系并不相同；两套数据的「同一电位」只有在覆盖度都达到稳态时才是同一个表面状态。

\subsubsection{测量条件与互补回查}

\begin{measure}
 \textbf{稳态与幅度}：体系在目标电位达到稳态后再扫频（判据见上一节测量条件第（5）条）；交流幅度取原书图 11-42／11-43 图注的 AC $10$ mV，并做第（2）条的幅度减半检验。
 \textbf{频率范围}：本分片读图实例所用原书图 11-42／11-43 的图注为 $20$ mHz 至 $5$ kHz；通用范围与点数的分配原则见上一节的测量条件第（3）条与规则三（源书 §11.5.1：低频端的点数要按「能否分辨出 45° 斜率」来定，而不是按对数均匀取点）。
 \textbf{电位区间}：至少覆盖目标反应从「无法拉第电流」到「扩散／钝化控制」的完整区间，才能用趋势论证速控步的转移（如读图实例中的 $0.10\sim0.48$ V）。
 \textbf{测量模式}：低阻抗、大电流体系（燃料电池单池、电池）宜用恒电流模式，避免恒电位仪过载；高阻抗体系（修饰电极、半导体电极）宜用恒电位模式\supp{恒电位／恒电流两种模式的适用判据属 §11.5.2 的讨论范围，不在本分片源文本内，此处只给判据}。
 \textbf{谱学联用前的采样条件}：若要与原位红外／拉曼的同一电位点对照，必须先让电位在该点稳定足够长时间（覆盖度达到稳态），否则两套数据的「同一电位」并不是同一个表面状态。
\end{measure}

\begin{crosscheck}
 阻抗谱的归属结论不应单独成立，至少与下列方法之一互查。上一节的互查清单（六条：$R_{\mathrm{ct}}$ 对 Tafel 斜率、扩散项对旋转圆盘变转速、$C_{\mathrm{dl}}$ 对循环伏安电容电流、$R_{\Omega}$ 对电流中断法、时间常数个数对独立状态变量个数、以及与其它分片的一致性）中，前两条已在上一节连同判据一起给出，此处不重复；真实数据读图另有表现的是三条：
 \textbf{（1）变温}：$R_{\mathrm{ct}}$ 的 Arrhenius 图给出表观活化能，可与同反应在稳态方法下的活化能比较。
 \textbf{（2）原位谱学}：第四象限的电感行为、低频吸附电容，应与原位红外／拉曼给出的 \ce{CO_{ad}} 覆盖度在同一电位区间出现——形状与分子信息对上了，归属才算落实。
 \textbf{（3）时间常数个数}：源文本强调，从阻抗谱中含有的时间常数个数及其数值大小可以推测电极过程的状态变量的情况，也可以观察电极过程中有无传质过程的影响，并由此得到溶液电阻、双电层电容及电极反应电阻等信息、推测反应机理\src{S2 §11.5，书内 p.866 / PDF p.354}（与上一节所引源书 p.856 处是同一条判据的两个出处）。把这条当作系统性的自检：时间常数个数应与你假定的机理中的独立状态变量个数一致，多了或少了都说明电路选错。
\end{crosscheck}

\subsection{电化学原位谱学技术}

常规电化学方法测的是界面上电流对电势或时间的响应，得到的是界面反应的综合信息，无法在分子水平上验证反应物、中间物种与产物，这给机理验证造成根本困难。原位谱学把谱学技术与电化学方法\textbf{同步}结合：在得到电流、电势等电化学信息的同时获得分子水平的信息，用于研究电极过程机理、电极表面特性、检测中间物种与产物、测定转移电子数、确定速率常数与扩散系数\src{S2 §11.6，书内 p.866—867 / PDF p.354—355}。本节按源文本的四个方面介绍：原位红外、原位拉曼、原位质谱与荧光单分子／单纳米粒子技术。

\begin{principle}
 原位（in situ）与离场（ex situ）的分界只有一条：\textbf{谱学测量是否在电化学调制进行的同时完成}。原位红外光谱电化学在反应体系被电化学调制与检测的同时监测红外信号变化；非原位红外则是在终止电化学反应之后对产物做红外研究\src{S2 §11.6.1，书内 p.867 / PDF p.355}。这一条分界决定了信息性质：原位方法能看到中间体与电位依赖，但必须解决「电解液对信号的干扰」与「传质受限」；离场方法信噪比与方法选择自由度大得多，但可能错过只在电位存在时才稳定的物种。
\end{principle}

\subsubsection{原位红外光谱电化学}

红外光谱电化学（infrared spectroscopy in electrochemistry）由红外光谱与电化学技术结合而成，可分为原位与非原位两类。原位红外既能检测电极表面吸附物种的身份与分子结构，也能探测电极表面及贴近表面薄液层中的反应物、产物与中间体，并为固／液界面提供双电层中物种的化学本质、成键方式及其随电位变化的信息。$1980$ 年 A. Bewick 等首次把电化学原位红外光谱成功用于电极／溶液界面研究，此后已建立的方法包括电化学调制红外光谱法（EMIRS）、差减归一化界面傅里叶变换红外光谱法（SNIFTIRS）、线性电位扫描反射光谱法（LPSIRS）、偏振调制红外光谱法（PMIRS）、傅里叶变换红外反射吸收光谱法（FTIRRAS）、时间分辨光谱电化学法（TRFTIRS）、二维红外光谱（2D-IR）与显微红外光谱电化学法（MFTIRS）等\src{S2 §11.6.1，书内 p.867 / PDF p.355}。

\begin{measure}
 \textbf{光路与窗口}：原位反射红外最常用外反射模式，红外辐射穿过窗口与溶液薄层，经研究电极表面反射后被检测。窗口材料必须对红外透明且不溶于所测溶液，如 \ce{CaF2}、\ce{Si}、\ce{ZnSe}。
 \textbf{薄层厚度是关键条件}：水及其他常用溶剂对红外辐射有强烈吸收，因此窗口与电极表面之间的液层必须很薄，约 $1\sim100$ nm。这一厚度同时决定了原位红外的时间响应与传质条件——薄层抑制了扩散，使本体浓度不再等于表面浓度，是原位红外定量结论的主要限制。
 \textbf{电极与池结构}：研究电极是封在玻璃或绝缘惰性材料中的圆盘，表面平整度要求高，固定在可调节电极与窗口距离的活塞末端，实验时充分接近窗口；通过调整反射镜角度选择最佳红外入射角，再用调制或差分技术取得最大检测信号\src{图 11-44、图 11-45，书内 p.868 / PDF p.356}。
\end{measure}

三种主要实验测试方式的差别在「减掉什么」上，理解这一点比记住缩写更重要\src{S2 §11.6.1，书内 p.868—870 / PDF p.356—358}：

\paragraph{EMIRS（电化学调制红外光谱）}
电势在「无研究物质产生」与「有电化学产物」两个位置之间以一定频率调制，用相敏检测器（PSD）解调，因而可以在测产物信号的同时排除溶剂与其他溶解物的吸收干扰，提高信噪比。测量的是电极电势由 $E_{1}$ 变为 $E_{2}$ 时表面相对反射率的改变量 $\Delta R$（相当于 $R(E_{2})-R(E_{1})$），最后以 $\Delta R/R$ 对波数作图。要点：\textbf{得到的是两个电位下的差谱}，因此单张谱不能直接给出绝对覆盖度。

\paragraph{SNIFTIRS（差减归一化界面傅里叶变换红外光谱）}
在傅里叶变换红外光谱基础上发展而来：单色仪被干涉仪取代，采集到的是强度随干涉仪物镜位置变化的干涉图；优点是可以多次采集并累加干涉图以提高信噪比。它测量两个不同电势下的谱图信号（电势分别取在所测物种的电化学过程发生与不发生处），在相应波长处作吸光度差减，再由计算机处理得到差谱。与 EMIRS 相比，它牺牲了调制解调的实时性，换来了可在宽波数范围同时取谱与更好的信噪比。

\paragraph{IRRAS（红外反射吸收光谱）}
依据红外反射光谱的表面选律：只有当吸附分子的偶极矩在电极表面法线方向上的分量不为零时，才能吸收平行于入射面的红外能量。装置上用光弹性调制器把入射光的偏振在 P 偏振与 S 偏振之间调制：P 光平行于入射面，可被电极表面与环境吸收；S 光垂直于入射面，不能与表面吸附物种相互作用，只包含环境对红外光的吸收。因此 P 光与 S 光的反射光谱之差就给出表面粒子的红外吸收信息。注意 IRRAS 在取谱时\textbf{电势保持不变}，得到的结果只代表该固定电势下的表面层信息，不同电位之间要各自成谱再比较。

\subsubsection{原位拉曼光谱}

\begin{principle}
 原位拉曼是在进行电化学过程的同时，用拉曼光谱研究电化学池中电极／溶液界面的性质：在测电位、电流获得电信号的同时，把一定波长的入射光照射在电极表面，再分析反射、透射或散射光束\src{S2 §11.6.2，书内 p.870 / PDF p.358}。
 散射的物理图像：光照射到分子或离子上，与分子发生碰撞并交换能量，光子的运动方向与能量随之改变。仅改变方向的弹性碰撞给出瑞利（Rayleigh）散射；方向与能量同时改变的非弹性碰撞给出拉曼（Raman）散射。散射光能量大于入射光能量的谱线（频率更高）称反斯托克斯线（紫伴线），小于入射光的称斯托克斯线（红伴线）\src{图 11-49，书内 p.871 / PDF p.359}。$1928$ 年 Raman 发现穿过透明介质的光频率发生改变，$1930$ 年获 Nobel 物理学奖。
 装置由拉曼光谱仪与原位电化学拉曼池两部分组成。光谱仪含激光源、收集系统（透镜组）、分光系统（光栅，或陷波滤光片结合光栅以滤除瑞利散射与杂散光）与检测系统（光电倍增管、半导体阵列检测器或多通道电荷耦合器件）。拉曼池具备工作电极、辅助电极、参比电极与通气装置，且必须配备光学窗口的密封体系，避免腐蚀性溶液与气体侵蚀仪器\src{S2 §11.6.2，书内 p.870—871 / PDF p.358—359}。
\end{principle}

\begin{measure}
 薄层条件是原位拉曼与原位红外共有的约束，但尺度不同：条件允许时应采用薄层溶液，电极与窗口间距取 $0.1\sim1$ mm，以尽量避免溶液信号的干扰；这一点对显微拉曼系统尤其重要，光学窗片或溶液层太厚会改变显微系统的光路，使表面拉曼信号的收集效率下降\src{S2 §11.6.2，书内 p.871 / PDF p.359}。
\end{measure}

\paragraph{应用一：CO 在贵金属表面的吸附}
CO 等小分子吸附在贵金属上时可视为偶极子，偶极子之间存在很强的相互作用（偶极耦合），导致 C--O 与 metal--CO 的拉曼频率蓝移；一般认为 metal--CO 的振动峰强于 C--O 峰。红外强度可代表金属上吸附 CO 的平均信号，而 SERS 研究发现了与红外类似的分压依赖行为，因此可以认为 SERS 信号能够代表纳米粒子上吸附 CO 分子的平均信号\src{S2 §11.6.2，书内 p.872 / PDF p.360}。这条结论是后面所有 SERS 定量讨论的前提，它的成立条件是「SERS 信号与分子数成正比」，而不是「SERS 峰强可比」。

\paragraph{应用二：酸性溶液中甲醇在 Pt 上的吸附与氧化}
在 $0.5$ mol/L \ce{CH3OH} + $0.5$ mol/L \ce{H2SO4} 溶液中扫甲醇氧化的循环伏安曲线并结合原位拉曼，可以得到甲醇在 Pt 上吸附氧化的完整图像\src{S2 §11.6.2，书内 p.872 / PDF p.360}：
\begin{obs}
 正扫至 $0.5$ V 之前，甲醇电氧化电流很小，可观察到 \qty{490}{cm^{-1}}、\qty{410}{cm^{-1}}、\qty{1060}{cm^{-1}} 及 \qty{1850}{cm^{-1}} 四个谱带；正扫过 $0.5$ V 后四个谱带明显减弱、氧化电流明显增大；扫至 $0.8$ V 时 CO 相关谱带几乎完全消失；扫至 $1.1$ V 时在 \qty{580}{cm^{-1}} 附近出现一个谱带。负扫时 \qty{580}{cm^{-1}} 谱带逐渐减弱，扫至 $0.8$ V 时消失，反扫的氧化电流出现峰值；负扫至 $0.5$ V 以下时 \qty{490}{cm^{-1}} 与 \qty{1060}{cm^{-1}} 谱带又逐渐增强。
\end{obs}
\begin{infer}
 四个谱带分别对应 Pt--CO$_{\mathrm{L}}$、Pt--CO$_{\mathrm{M}}$、C--O$_{\mathrm{L}}$ 与 C--O$_{\mathrm{M}}$ 的伸缩振动，说明该电位下电极表面 Pt 处于 CO「中毒」状态；谱带减弱与电流增大同步，说明电势扫至 $0.5$ V 时表面上电化学生成的 CO 开始被氧化；$0.8$ V 时 CO 谱带消失说明 CO 已完全氧化。$1.1$ V 出现的 \qty{580}{cm^{-1}} 谱带来源于 Pt 的多种氧化物种的 Pt--O 振动；负扫时它减弱并消失，说明氧化的 Pt 被还原、活性位增加，因此反扫氧化电流出现峰值。负扫至 $0.5$ V 以下时 Pt--CO 与 C--O 谱带重新增强，说明低电势下甲醇重新吸附于电极表面、脱氢后重新生成 CO；而 CO 在较低电势下无法被氧化成 \ce{CO2}，于是电极被 CO 毒化，阻碍甲醇进一步脱氢，氧化电流下降。
\end{infer}
这里\textbf{观察与推断的接口是谱带归属}：四个谱带（\qty{490}{cm^{-1}}、\qty{410}{cm^{-1}}、\qty{1060}{cm^{-1}}、\qty{1850}{cm^{-1}}）与 Pt--CO$_{\mathrm{L}}$、Pt--CO$_{\mathrm{M}}$、C--O$_{\mathrm{L}}$、C--O$_{\mathrm{M}}$ 伸缩振动的对应关系来自文献与标定实验，一旦归属错位，整段机理推断都会跟着错。因此这类工作的可信度取决于是否在同一体系里做了归属的独立验证（如同位素取代、参比化合物或理论计算）。

\begin{pitfall}
 \textbf{（1）拉曼信号弱是方法的固有短板。} 拉曼散射光子数极低——源文本给出的量级是荧光的 \num{1e-14}、红外的 \num{1e-6}，如此低强度的散射光用于研究表面或界面分子的结构与性质十分困难\src{S2 §11.6.2，书内 p.872—873 / PDF p.360—361}。常用对策是表面增强拉曼光谱（SERS）：把试样吸附于胶态金属（Cu、Ag、Au）或这些金属的粗糙表面上，再用普通拉曼方法测量，强度可达正常情况的 \num{1e5}～\num{1e6} 倍，因而可用于极低浓度的样品溶液，并能有效避免溶液中相同物种的信号干扰。
 \textbf{（2）SERS 的增强是条件依赖的，不是常数。} 增强因子取决于基底形貌、金属种类、激发波长与吸附构型，不同批次基底之间不可直接比较峰强。用 SERS 峰强做定量时，必须同批基底、原位标定；把 \num{1e5}～\num{1e6} 当作一个固定倍数去反算表面浓度是最常见的误用。
 \textbf{（3）薄层条件本身会改变被测体系。} 电极与窗口间距 $0.1\sim1$ mm 的薄层抑制了自然对流与本体扩散，使电极表面浓度低于本体；因此薄层谱学给出的电位—覆盖度关系不能直接外推到流动或厚层条件下的动力学。
 \textbf{（4）研究条件与实际工作条件差别太大。} 燃料电池基础研究一般在室温下用三电极体系，而组装成电堆后为提高反应速率常采用高温运行；考虑到 Nafion 膜的稳定性与活性以及三相水界面的要求，质子交换膜燃料电池的使用温度一般在 $80$ \degC 左右。温度不仅按 Arrhenius 关系改变反应速率，也会改变反应途径——甲醇、甲酸的直接途径与经 CO 的间接途径，氧还原的 2 电子与 4 电子过程，不同途径给出的电化学性能不同。若研究条件与实际情况相差太大，就很难由研究条件下的表征结果推断实际工作条件下的行为，因此有必要研制变温流动光谱电解池来模拟燃料电池的实际工作状态\src{S2 §11.6.2，书内 p.873 / PDF p.361}。
\end{pitfall}

\subsubsection{原位质谱}

质谱与电化学反应装置联用最早出现在 $20$ 世纪 $70$ 年代。$1971$ 年 Bruckenstein 和 Gadde 用气相质谱检测电化学反应产生的挥发性物质，率先建立电化学质谱技术（EMS）：把 Pt 工作电极制备在多孔玻璃膜上，采用膜进样方式，用多孔 Teflon 膜把电解液与质谱真空进样系统分隔开，只允许挥发性气体产物进入电离室；其特点是气体样品累积到一定量后一次性导入质谱作定性与定量分析，测量电极上气体产物的响应时间约 $20$ s。$1984$ 年 Wolter 和 Heitbaum 建立微分电化学质谱（DEMS），沿用多孔 Teflon 膜进样，区别在于用两级真空泵分别对电离室与质量分析系统分级抽真空，进样口利用压差把反应产生的气体快速吸入电离室，得到的质谱信号强度正比于电化学反应的法拉第电流；DEMS 实现了对挥发性产物的\textbf{连续}测量，总响应时间小于 $1$ s\src{S2 §11.6.3，书内 p.873—874 / PDF p.361—362}。

\begin{principle}
 DEMS 系统由电化学反应装置、膜接口与质谱仪三部分组成，膜接口是核心设计，实现液相与真空的连接：产物在膜两侧压差作用下从液相挥发进入真空系统，采样耗时极短，因而可在线检测。其定量基础是「离子电流正比于法拉第电流」，推导链条为\src{原书式 11-126—11-129，书内 p.874 / PDF p.362}：
 \begin{align}
  I_{\mathrm{F}}&=\frac{\mathrm{d}Q_{\mathrm{F}}}{\mathrm{d}t}=\frac{\mathrm{d}(zFn)}{\mathrm{d}t}=zF\frac{\mathrm{d}n}{\mathrm{d}t}
  \label{eq:dems-if}\\
  n_{\mathrm{i}}&=K_{1}K_{2}n
  \label{eq:dems-ni}\\
  Q_{\mathrm{i}}&=K_{3}n_{\mathrm{i}}
  \label{eq:dems-qi}\\
  I_{\mathrm{i}}&=\frac{\mathrm{d}Q_{\mathrm{i}}}{\mathrm{d}t}
   =K_{1}K_{2}K_{3}\frac{\mathrm{d}n}{\mathrm{d}t}
   =\frac{K_{1}K_{2}K_{3}}{zF}I_{\mathrm{F}}
  \label{eq:dems-ii}
 \end{align}
 令 $K=K_{1}K_{2}K_{3}/(zF)$，即得 $I_{\mathrm{i}}=KI_{\mathrm{F}}$；若法拉第电流效率不是 $100\%$，电流效率常数也应计入系数。式中 $K_{1}$ 为实际被采样的比例系数（产物在压差下进入真空系统的同时也会向溶液本体扩散），$K_{2}$ 概括「进入离子源」与「被电离进入质量分析器」两级比例。\tocheck{未对照原页面图像：源文本在 $K_{1}$ 与 $K_{2}$ 之间的定义句有脱字（OCR 作「可将这两项系数统一考虑，定义由此可得到」），$K_{2}$ 的严格定义需对照原书 p.874}
 由于影响 $K$ 值的因素众多（电化学实验环境、质谱仪参数、物种自身性质），不同物种在不同环境下的 $K$ 值可能都不相同，因此\textbf{特定物种的 $K$ 值必须在电化学条件下定量生成该物种、由质谱检测其离子电流来标定}，不能沿用文献值\src{S2 §11.6.3，书内 p.874 / PDF p.362}。
\end{principle}

对非挥发性产物的在线检测，$1986$ 年 Hambitzer 和 Heitbaum 把 N,\,N-二甲基苯胺电解氧化，并将电极表面的电解液连续导入质谱热喷雾离子室的前置进样毛细管中，实现了在线检测。随着质谱技术发展，电化学质谱联用结合了更多样品电离方式以分析不同形态与性质的产物，包括电喷雾电离、大气压化学电离、大气压光致电离、快原子轰击电离与电感耦合等离子体电离等。另一种可与更常用电极材料匹配的设计，是把 Teflon 膜接口放在距旋转圆柱电极很近处（约 $0.3$ mm）\src{S2 §11.6.3，书内 p.874—875 / PDF p.362—363}。

二次离子质谱（SIMS）是表征表面与薄膜的非原位 UHV 方法，可与电化学联用：用高能一次离子束（如 $15$ keV \ce{Cs+}）轰击表面，通过溅射侵蚀表面、由表面组成中产生二次离子并用质谱检测。利用一次离子束扫描可作二维表征，监测单离子强度与溅射时间的关系可得深度分布。SIMS 的测量极限比 XPS 或 AES 高得多（\num{1e-8}～\num{1e-4} 原子分数）。但 SIMS 并不是一种「真实表面」技术：二次离子的产生效率由一次离子束所产生的薄的离子注入层的三维性质决定，深度分布的假象可由此产生。把电极转移到 UHV 样品腔后，研究电极表面的其他可用方法是热脱附与激光脱附\src{S2 §11.6.3，书内 p.875 / PDF p.363}。

\begin{crosscheck}
 原位质谱与谱学方法的分工很清楚：红外／拉曼给\textbf{表面物种与成键}，质谱给\textbf{产物种类与生成速率}。二者对同一电位点必须给出一致的速率控制图像——若 DEMS 显示某产物在 $0.5$ V 起快速生成，而原位红外显示中间体覆盖度在同一电位开始下降，则两者共同支持「中间体在该电位被消耗」的机理。反过来，如果只报 DEMS 的产物分布而不报 $K$ 的标定方式，产物比例（常被用来判断 2 电子与 4 电子途径的比例）就只是相对信号，不能作为机理证据。
\end{crosscheck}

\subsubsection{荧光单分子与单纳米粒子催化}

纳米材料与催化领域的发展，需要能在微观层面解析纳米催化剂催化过程的新技术。荧光单分子／单纳米粒子技术基于荧光显微技术，利用现代单光子检测器的高灵敏度，实现单个纳米粒子表面荧光产物分子一个一个的观测，从而从微观反应动力学与热力学角度解析催化过程\src{S2 §11.6.4，书内 p.875—876 / PDF p.363—364}。

\begin{principle}
 该技术跟踪单个分子的荧光信号：荧光相关的前驱体分子在被催化发生氧化还原反应的过程中，在荧光态与非荧光态之间转化，通过跟踪荧光信号的强弱获取反应动力学信息。要拿到单分子与单纳米粒子的信号，需要同时满足三个条件\src{S2 §11.6.4，书内 p.877 / PDF p.365}：
 （1）\textbf{浓度极低}——催化剂浓度与反应物浓度都要极低，确保每次捕获到的是单次反应引起的荧光信号变化；
 （2）\textbf{界面全反射激发}——利用激光在界面处发生全反射时产生的渐逝波（倏逝波）激发界面处的荧光分子，避免背景溶液中的荧光分子被激发而抬高背景；
 （3）\textbf{超快单光子检测}——用现代单光子检测器捕捉信号，在获得良好信噪比之后，对大量荧光信号作统计分析，得到产物形成与脱附过程的信息。
 具体做法是把纳米粒子稀疏分散在石英载玻片或类似基底上，反应物流过表面时在纳米粒子表面活性位上生成可被激光激发荧光的产物分子。产物生成之前看不到该纳米粒子；催化反应一旦完成，生成的荧光分子立刻被激发发光，因而可以通过单分子荧光显微镜直接观察检测，并通过定位该荧光分子精确确定该纳米粒子的位置。用单光子检测器可实时记录单个产物分子\textbf{形成前的等待时间}与\textbf{形成后到脱附前在表面的停留时间}——这样就把常规方法中的「一个催化周期」拆成两个更细致的过程：产物分子的生成，以及产物的脱附与活性位的释放，从而分别考察单个纳米粒子表面的产物形成速率与脱附速率\src{S2 §11.6.4，书内 p.877 / PDF p.365}。
\end{principle}

最典型的系统研究出现在 $2008$ 年：陈鹏等以金纳米粒子为催化模型，首次用荧光单分子催化技术在单分子水平上研究了单个金纳米粒子催化羟胺与刃天青（resazurin）反应生成荧光产物试卤灵（resorufin）的过程，并首次提出适用于单分子单纳米粒子催化的动力学机理模型\src{S2 §11.6.4，图 11-51，书内 p.876—877 / PDF p.364—365}。该方法的独特价值在于能一个一个地研究看似一样的纳米粒子，鉴别粒子之间的活性差别（反应快慢、产物吸附强度或脱附速率），而这些信息是常规方法无法获得的；除了用统计方法得到大量粒子的平均信息，还能得出个别粒子的「个性」信息，有助于更精确地理解催化剂特性与其物理特性（粒径、组成、晶面类型）之间的相关性\src{S2 §11.6.4，书内 p.878 / PDF p.366}。

在超分辨成像方向，$2006$ 年 Hofkens 等用荧光显微镜实时观察催化反应过程中单个有机分子的化学转化，绘制出催化剂晶体的活性位点分布，并发现不同催化晶面具有催化不同化学反应的能力；$2009$ 年该组进一步把超分辨成像精度提高到可与电子显微镜匹敌；$2013$ 年 Majima 课题组用自主合成的探针对 \ce{TiO2} 光解水过程的活性位点作超分辨成像，发现可见光引起的还原反应发生在金纳米粒子附近几十纳米范围内，在微观层面证实了等离子引发的电子空穴分离机理；$2014$ 年方宁课题组用荧光超分辨成像揭示单个金—硫化镉纳米粒子表面的光催化性质，揭示了两种完全不同的、能量依赖的电荷分离机理；$2015$ 年 Hofkens 等对石油裂解过程产物分子作超分辨成像，发现分子筛催化石油裂解过程中产物主要在布朗斯特酸位点上形成，并揭示了不同粒子间的差异性。超分辨成像不仅能描绘静态活性位点，还能跟踪其位置变化，例如中国科学院长春应用化学研究所徐维林课题组 $2015$ 年对纳米金三角等催化剂纳米粒子作动态超分辨成像，发现金纳米粒子表面活性位点因重构而发生移动\src{S2 §11.6.4，书内 p.878 / PDF p.366}。

在动力学方向，$2008$ 年陈鹏等利用催化剂在荧光态与非荧光态之间转化的过程，实现对微观催化反应的动力学分析并揭示了三种不同的催化机理；$2014$ 年徐维林课题组首次结合量化计算与「用空心纳米笼困住反应分子」的方法，实现化学反应过程中各基元步骤的分子反应动力学分析；随后该组把荧光单分子催化技术拓展到电化学领域，揭示了电催化剂的衰减机理；同年该课题组通过对不同粒径纳米立方钯催化剂催化活性的解析，获得纳米粒子「边—面」活性差异，把单粒子催化技术拓展到亚单粒子水平。该技术虽刚发展起来，其应用正朝着更深更广的方向发展，已从催化动力学研究扩展到热力学研究\src{S2 §11.6.4，书内 p.879 / PDF p.367}。

\subsubsection{原位谱学的共同边界}

\begin{table}[htbp]
 \centering
 \caption{四类原位谱学方法的分工与边界。「关键条件」一列是使用该方法数据前必须确认的前提；「主要陷阱」一列给出最常见的误用。\src{据 S2 §11.6.1—§11.6.4 汇总，书内 p.867—879 / PDF p.355—367。}}
 \label{tab:ch11c-insitu-methods}
 \begin{tabularx}{\onecolwidth}{@{}p{20mm}p{29mm}p{30mm}X@{}}
  \toprule
  方法 & 拿到什么 & 关键条件 & 主要陷阱 \\
  \midrule
  原位红外（EMIRS／SNIFTIRS／IRRAS） & 吸附物种身份与成键、双电层物种随电位的变化 & 窗口材料透明且不溶；薄层 $1\sim100$ nm & 得到的是两个电位的差谱（EMIRS／SNIFTIRS），不能直接读绝对覆盖度；IRRAS 只代表固定电位 \\
  原位拉曼（含 SERS） & 振动指纹：金属—吸附物与分子内振动 & 薄层 $0.1\sim1$ mm；SERS 需粗糙／胶态贵金属基底 & 增强因子随基底与波长变化，峰强不可跨批比较；薄层改变传质 \\
  原位质谱（EMS／DEMS） & 挥发性产物的种类与生成速率；可作同位素标记 & 膜接口与压差；响应时间 $<1$ s（DEMS） & 系数 $K$ 必须原位标定；只测挥发性物种，非挥发性产物需另接电离方式 \\
  SIMS 与脱附法（非原位 UHV） & 表面组成、深度分布、脱附物种 & 需把电极转移进 UHV & SIMS 由注入层三维性质决定产额，会产生深度分布假象；不是真实表面技术 \\
  荧光单分子／单粒子 & 单粒子层面的生成与脱附动力学、活性位点空间分布 & 极低浓度；全反射倏逝波激发；单光子检测 & 分析对象是荧光产物分子，只适用于能生成荧光产物的反应体系 \\
  \bottomrule
 \end{tabularx}
\end{table}

\begin{pitfall}
 原位谱学的共同陷阱不是「分辨率不够」，而是\textbf{为了能测而改变了被测对象}：红外要求纳米级薄层、拉曼要求亚毫米级薄层、质谱要求真空接口、荧光要求极低浓度，每一项都在改变传质与表面状态。因此原位谱学给出的电位—覆盖度或电位—速率关系，只在同样的薄层与浓度条件下成立；把它们直接外推到厚层、流动或高电流工况（这正是燃料电池的真实工作状态），是这一类方法最容易被忽略的边界\src{S2 §11.6.2，书内 p.873 / PDF p.361}。
\end{pitfall}

\begin{crosscheck}
 原位谱学的结论应与三类独立信息互查：\textbf{（1）电化学响应}——谱带出现／消失的电位应与循环伏安上的电流峰或阶跃位置一致（如甲醇氧化中 $0.5$ V 的谱带减弱与电流增大同步）；\textbf{（2）非原位表征}——把同一电极取出做 XPS、SIMS 或脱附实验，检查表面物种与原位结论是否一致，并注意 UHV 转移过程中的表面变化；\textbf{（3）同位素标记}——用 \ce{^{13}C} 或 \ce{D} 取代确认谱带归属与产物来源，这是把「谱带—物种」对应关系从文献搬运变成自证的关键一步。
\end{crosscheck}

\srcnote{本分片覆盖 S2（下册）第 11 章 §11.5 后半（阻抗谱的应用与读图，书内 p.866 / PDF p.354）与 §11.6 全节（原位谱学技术及应用，书内 p.866—879 / PDF p.354—367）。§11.5.1 的阻抗基本定义、等效电路推导与 Nyquist／Bode 图定义不在本分片源文本范围内——它们写在本章另一分片（§11.5.1，依源文本写成、逐条标有 \src{} 依据），故本分片不再重复，只保留源文本 §11.5.2 及以后特有的内容：读图实例（甲醇氧化）、参数相关性与误判、测量条件与互补回查、以及 §11.6 全节；其中图 11-42／11-43 的坐标与相角取值、以及 DEMS 系数 $K_{1}$、$K_{2}$ 的定义按学科标准形式重建，存疑点已在正文用 \tocheck{} 就地标出。图 11-42、11-43、11-44～11-51 未对照原页面图像，本分片不再自绘阻抗模型图（原型计算的 Nyquist 模型图见上一节），未裁切原书谱图。}
